% Source: polynomial-PEPS manuscript, Proposition 4.1, 03-patches.tex,
% lines 563--603, revision adc7f1241b42e322a6451854ab7e4b4c146bf78a.
\begin{theorem}[Square patches from nested cylinder spaces]
    \label{thm:peps_nested_cylinder_square_patch}
    \leanok
    \lean{TNLean.PEPS.EncodedFrame.exists_squarePatch_of_mem_iSup_dependentRegionCylinder}
    Let $V$ be a finite set of sites with positions $p_v\in\mathbb R^2$, and let
    $\mathcal H_R=\mathbb C^{R\to\{0,\ldots,q-1\}}$ for $R\subseteq V$ and $q\in\mathbb N$.
    Fix a centre $z$, a scale $u>0$, and $m\in\mathbb N$.
    Let $(r_j)_{0\leq j<m}$ be nondecreasing, with $u\leq r_j\leq2u$ for each $j$.
    Put $R_j=\{v\in V:\|p_v-z\|_\infty\leq r_j\}$, and let
    $S_j\subseteq\mathcal H_{R_j}$ be subspaces.
    Suppose $\Omega,v\in\mathcal H_V$ and
    $v\in\sum_{j<m}\mathcal C_{R_j}(S_j)$.
    Then there is a square patch with centre $z$, scale $u$, exactly these
    $m$ radii, and inside dimensions $d_j$ such that
    \begin{align}
        \sum_{j<m}d_j   & \leq\sum_{j<m}\dim S_j, \notag \\
        \|(I-Q)\Omega\| & \leq\|\Omega-v\|, \notag
    \end{align}
    where its cylinder sum is
    \begin{align}
        Q & =\sum_{j<m}\sum_{\ell<d_j}
        \bigl(\lvert v_{j\ell}\rangle\langle v_{j\ell}\rvert\bigr)_{R_j}
        \otimes I_{V\setminus R_j}. \notag
    \end{align}
    The sum is zero when the family is empty.
\end{theorem}
\begin{proof}\leanok
    \uses{thm:peps_nested_cylinder_orthogonalization,
        thm:peps_nested_cylinder_projector_sum,
        thm:peps_coordinate_reset_gram,thm:peps_cylinder_projector}
    Replace the inside spaces by their orthogonalized projected images
    $\widetilde S_j$, and choose an orthonormal basis in each image.
    Their cylinders are mutually orthogonal and have the same span
    $\mathcal E=\sum_{j<m}\mathcal C_{R_j}(S_j)$.
    Thus the lifted rank-one projectors are pairwise orthogonal, and their
    sum is $Q=P_{\mathcal E}$.
    The bounds $\dim\widetilde S_j\leq\dim S_j$ give the dimension estimate.
    Since $v\in\mathcal E$, orthogonal projection onto $\mathcal E^\perp$ gives
    \begin{align}
        \|(I-Q)\Omega\|
         & =\|P_{\mathcal E^\perp}(\Omega-v)\|
        \leq\|\Omega-v\|. \notag
    \end{align}
    The canonical identification of configurations preserves these operators
    and the Hilbert-space norm.
\end{proof}
